\section{总结与延伸}

\subsection{讲者的核心总结}

\begin{figure}[H]
\centering
\includegraphics[width=0.95\textwidth]{slides-images/slide-051.jpg}
\caption{全课总结\protect\footnotemark}
\end{figure}
\footnotetext{Slides 第52页。}

Tatsu Hashimoto 在本讲中传达了以下核心信息：

\begin{enumerate}
    \item \textbf{数据 Scaling 有坚实的理论基础}：幂律衰减可以从非参数统计的基本定理中推导出来，缩放指数与数据的内在维度相关
    \item \textbf{模型/架构 Scaling 提供了强大的工程工具}：通过在小规模上比较不同选择，可以可靠地预测大规模下的最优方案
    \item \textbf{Chinchilla 确立了资源分配的基本框架}：联合数据--模型 Scaling Laws 回答了``给定预算，模型应该多大''的核心问题
    \item \textbf{实践已超越 Chinchilla}：推理成本的考量使得业界大幅``过度训练''小模型，tokens/parameter 比例从 20 增长到数百甚至数千
\end{enumerate}

\subsection{全课知识图谱}

\begin{center}
\begin{tikzpicture}[
    node distance=1.8cm and 2.5cm,
    every node/.style={draw, rounded corners, align=center, minimum width=2.8cm, minimum height=0.8cm, font=\small},
    arrow/.style={->, thick}
]
    \node (theory) {统计学习理论\\（VC维、非参数率）};
    \node (datascale) [below=of theory] {数据 Scaling\\$L \propto D^{-\alpha}$};
    \node (modelscale) [right=of datascale] {模型 Scaling\\架构/超参比较};
    \node (joint) [below=of datascale, xshift=2.2cm] {联合 Scaling\\Rosenfeld/Kaplan};
    \node (chinchilla) [below=of joint] {Chinchilla\\20 tokens/param};
    \node (inference) [right=of chinchilla] {推理最优\\100+ tokens/param};
    \node (practice) [below=of chinchilla, xshift=2.2cm] {工程决策流程\\小模型 $\to$ 外推};

    \draw[arrow] (theory) -- (datascale);
    \draw[arrow] (datascale) -- (joint);
    \draw[arrow] (modelscale) -- (joint);
    \draw[arrow] (joint) -- (chinchilla);
    \draw[arrow] (chinchilla) -- (inference);
    \draw[arrow] (chinchilla) -- (practice);
    \draw[arrow] (inference) -- (practice);
\end{tikzpicture}
\end{center}

\subsection{拓展阅读}

\begin{itemize}
    \item Kaplan et al., \textit{Scaling Laws for Neural Language Models}: \url{https://arxiv.org/abs/2001.08361} —— Scaling Laws 领域最全面的参考论文
    \item Hoffmann et al., \textit{Training Compute-Optimal Large Language Models} (Chinchilla): \url{https://arxiv.org/abs/2203.15556} —— 确立 20 tokens/parameter 比例的里程碑论文
    \item Hestness et al., \textit{Deep Learning Scaling is Predictable, Empirically} (2017): \url{https://arxiv.org/abs/1712.00409} —— 大规模神经网络 Scaling Law 的早期系统研究
    \item Rosenfeld et al., \textit{A Constructive Prediction of the Generalization Error Across Scales} (2020): \url{https://arxiv.org/abs/1909.12673} —— 联合 Scaling Law 函数形式的重要参考
    \item Epoch AI, \textit{Chinchilla Scaling: A Replication Attempt} (2024): \url{https://epochai.org/blog/chinchilla-scaling-a-replication-attempt} —— 对 Chinchilla 方法 3 的复现与修正
    \item Yang et al., \textit{Tensor Programs V: Tuning Large Neural Networks via Zero-Shot Hyperparameter Transfer} (MuP): \url{https://arxiv.org/abs/2203.03466} —— 跨尺度学习率迁移的理论与方法
    \item Muennighoff et al., \textit{Scaling Data-Constrained Language Models}: \url{https://arxiv.org/abs/2305.16264} —— 多 Epoch 训练的 Scaling Laws
    \item Banko \& Brill, \textit{Scaling to Very Very Large Corpora for Natural Language Disambiguation} (2001) —— 数据缩放的早期 NLP 研究
\end{itemize}

\end{document}
